\section{Characteristic-class constraints in pseudoindex five}
\label{sec:characteristic-classes}

We now assume that $X$ is a Fano sixfold with nef tangent bundle,
$\rho(X)=1$, and pseudoindex five. Let
\[
 \pi:U\longrightarrow V,\qquad e:U\longrightarrow X
\]
be its smooth universal family. Thus $\dim U=9$, $\dim V=8$,
the fibres of $\pi$ are projective lines, and $e$ has relative dimension
three. The preceding geometric argument gives
$T_e|_C\simeq\mathcal O_{\mathbb P^1}(-1)^{\oplus3}$ for every
$\pi$-fibre $C$. Our goal is to constrain the polarization of an
evaluation fibre by combining the two fibrations.

Put
\begin{equation}\label{eq:d5-generators}
 h_X=\frac{c_1(T_X)}5,\qquad h=e^*h_X,\qquad
 c=c_1(T_\pi),\qquad \ell=2h-c,\qquad E=T_e.
\end{equation}
The class $h$ has degree one on a $\pi$-fibre. Lemma~\ref{lem:p1-involution}
therefore shows that $\ell$ comes from $V$ and that $\bar h=\ell-h$.

\subsection{The polarized evaluation fibre}

\begin{lemma}\label{lem:marked-tangent}
For a general point $x\in X$, the fibre $F=e^{-1}(x)$ is a smooth
irreducible projective threefold, and the tangent direction at the marking
defines a finite morphism
\[
 \tau_x\colon F\longrightarrow\mathbb P(T_x^*X).
\]
Its image $\mathcal C_x$ is the variety of minimal rational tangents of
$V$ at $x$.  The line bundle $L=\tau_x^*\mathcal O(1)$ is ample and
generated by global sections, with $c_1(L)=\ell|_F$.
\end{lemma}

\begin{proof}
Smoothness, projectivity and connectedness follow from the established
properties of $e$.  A smooth connected variety is irreducible.
Let $V_x\subset V$ be the reduced locus of curves passing through $x$.
The restriction $\pi|_F\colon F\to V_x$ is proper and has finite fibres:
each curve has only finitely many points on its normalization mapping to
$x$.  It is therefore finite.  Kebekus's Theorem~3.3(1)
\cite{Kebekus2002}, applied with the ample line bundle
$\omega_X^{-1}$, says that only finitely many members of $V_x$ can be
singular at a general $x$, and that even these curves are immersed at
their markings.  Since $\dim V_x=3$, a general member is smooth at $x$,
so $\pi|_F$ is birational.  As $F$ is normal, it is the normalization of
$V_x$.  The hypotheses of Kebekus's Theorem~3.4 are satisfied: $V$ is an
irreducible covering family and $V_x$ is proper.  That theorem gives the
finite tangent morphism on $F$.

With our quotient convention for projectivization,
$\mathbb P(T_x^*X)$ parametrizes lines in $T_xX$.  The differential of
$e$ is nonzero on $T_\pi$ at every marking over $x$, and identifies
$T_\pi|_F$ with the pullback of the tautological subline
$\mathcal O(-1)\subset T_xX\otimes\mathcal O$.  Consequently
$c_1(L)=-c|_F=\ell|_F$.  Pullback by a finite morphism preserves
ampleness, and the pullbacks of hyperplane sections generate $L$.
\end{proof}

\subsection{Classes descending to the parameter space}

Define $r,s,t$ to be the elementary symmetric polynomials in the three
formal roots obtained by adding $h$ to the Chern roots of $E$.
Equivalently,
\begin{align*}
 r&=c_1(E)+3h,\\
 s&=c_2(E)+2hc_1(E)+3h^2,\\
 t&=c_3(E)+hc_2(E)+h^2c_1(E)+h^3.
\end{align*}
This definition uses only rational cohomology; it does not assert that
$h$ is the first Chern class of a line bundle.

\begin{lemma}\label{lem:shifted-descent}
The classes $r,s,t$ are pulled back from $V$.
\end{lemma}
\begin{proof}
There is an actual line bundle
\[
 \mathcal A=e^*\omega_X^{-1}\otimes(T_\pi^\vee)^{\otimes2},
 \qquad c_1(\mathcal A)=5h-2c=h+2\ell.
\]
Its degree on every $\pi$-fibre is $5-4=1$.
Thus $E\otimes\mathcal A$ restricts to the trivial rank-three bundle
on each fibre. Cohomology and base change show that
$\mathcal B=\pi_*(E\otimes\mathcal A)$ is locally free of rank three.
The evaluation map $\pi^*\mathcal B\to E\otimes\mathcal A$ is an
isomorphism on every fibre, and hence is an isomorphism.
The formal roots defining $r,s,t$ are consequently the roots of
$\pi^*\mathcal B$ minus $2\ell$. Their elementary symmetric
polynomials are pulled back from $V$.
\end{proof}

For the fibre $F$ and polarization $L$ of
Lemma~\ref{lem:marked-tangent}, the definitions and $h|_F=0$ give
\begin{equation}\label{eq:fibre-polarization}
 c_1(L)=\ell|_F,
 \qquad r|_F=c_1(T_F).
\end{equation}
We shall constrain the ratio
\begin{equation}\label{eq:ratio-u}
 u=\frac{\int_F c_1(T_F)c_1(L)^2}{\int_F c_1(L)^3}.
\end{equation}
Its denominator is positive.
The intersection identities below will force a quadratic equation for
$u$. In Section~\ref{sec:grassmannian}, this equation and the
sectional-genus formula will determine the degree of $L$.

\subsection{Two direct-image identities}

We first record why the Hodge-theoretic direct images appearing in
Riemann--Roch contribute no positive-degree Chern character.

\begin{lemma}\label{lem:constant-hodge}
Let $f:Z\to B$ be a smooth projective morphism between smooth connected
complex projective varieties, and suppose that $B$ is simply connected.
For all $p,j\ge0$, the vector bundle $R^jf_*\Omega^p_{Z/B}$ is
holomorphically trivial. In particular its positive-degree rational
Chern character vanishes.
\end{lemma}
\begin{proof}
The local systems $R^kf_*\mathbb Q$ are constant. The global invariant
cycle theorem gives a surjection
\[
 H^k(Z,\mathbb Q)\longrightarrow H^0(B,R^kf_*\mathbb Q)
   \simeq H^k(Z_b,\mathbb Q)
\]
for every $b\in B$; see Deligne~\cite[Theorem~4.1.1(ii)]{Deligne1971}.
Restriction is a morphism of pure Hodge structures. Therefore each
Hodge component of $H^k(Z_b,\mathbb C)$ is spanned by restrictions of
classes of the same type on $Z$. A class on $Z$ restricts to a flat
section of the local system, and a class of fixed Hodge type keeps that
type on every fibre. Thus the Hodge decomposition, and hence the Hodge
filtration, is constant in the flat trivialization. This is also the
fixed-part statement of~\cite[Corollary~4.1.2]{Deligne1971}.
The relative Hodge--de Rham degeneration identifies
$R^jf_*\Omega^p_{Z/B}$ with the degree-$p$ graded bundle of the Hodge
filtration on $R^{p+j}f_*\mathbb C\otimes\mathcal O_B$; see
\cite[Section~2(a)--(c)]{Griffiths1970}. Since that
filtration is constant, each graded bundle is trivial.
\end{proof}

The base $X$ is simply connected, so the lemma applies to $e$.
Relative Grothendieck--Riemann--Roch
\cite[\S7, p.~113]{BorelSerre1958} for $\mathcal O_U$ and
$\Omega^1_{U/X}=E^\vee$ gives
\begin{equation}\label{eq:grr-vanishings}
 e_*\operatorname{td}_i(E)=0,
 \qquad
 e_*[\operatorname{td}(E)\operatorname{ch}(E^\vee)]_i=0
 \quad(i>3).
\end{equation}
Indeed these are the components of weight $i-3>0$ of
\[
 \operatorname{ch}\!\left(\sum_j(-1)^jR^je_*\mathcal O_U\right)
 \quad\text{and}\quad
 \operatorname{ch}\!\left(\sum_j(-1)^jR^je_*\Omega^1_{U/X}\right),
\]
respectively.

These two choices capture the equations from all relative holomorphic
forms. Indeed, for a rank-three bundle and $p=0,1$,
$\bigwedge^{3-p}E^\vee=\det(E^\vee)\otimes\bigwedge^pE$ and
$\operatorname{td}(E^\vee)=e^{-c_1(E)}\operatorname{td}(E)$.
Replacing the Chern roots by their negatives therefore gives
\[
 [\operatorname{td}(E)\operatorname{ch}(\textstyle\bigwedge^{3-p}E^\vee)]_i
 =(-1)^i
 [\operatorname{td}(E)\operatorname{ch}(\textstyle\bigwedge^pE^\vee)]_i.
\]
Thus $\Omega^2_{U/X}$ and $\Omega^3_{U/X}$ repeat, up to sign in
each weight, the equations from $\Omega^1_{U/X}$ and $\mathcal O_U$.

\subsection{Explicit polynomial representatives}

We now turn the tangent-sequence and Riemann--Roch identities into
linear relations among weight-nine intersection numbers on $U$. The
projection formula will then identify selected normalized numbers as
products involving $u$.
Set
\[
 x_i=e^*(i!\operatorname{ch}_i(T_X)),\qquad i\ge1.
\]
Then $x_1=5h$ and $x_i=0$ for $i\ge7$, by dimension. Work in the
polynomial ring
\[
 R=\mathbb Q[h,x_2,x_3,x_4,x_5,x_6,\ell,r,s,t],
\]
with respective weights $(1,2,3,4,5,6,1,1,2,3)$.
Evaluation of these generators as cohomology classes defines a graded
homomorphism $R\to H^{2*}(U,\mathbb Q)$.

Our unknowns are the integrals of weight-nine monomials in $R$.
We obtain linear relations among them from three sources: classes
descending to the eight-dimensional space $V$, classes pulled back
from the six-dimensional space $X$ and their $\pi$-integrals, and the
two Riemann--Roch vanishings. The tangent sequences make the
$\pi$-integrals explicit. Recall that
$\delta=\pi^*\pi_*$ lowers weight by one. The recurrences below
express these operations as polynomials, before we collect the
resulting relations into a matrix.

Let $p_j$ be the Newton power sum of the three formal roots with
elementary symmetric functions $r,s,t$, and put $p_0=3$.
For complete specificity, if $s_1=r,s_2=s,s_3=t$, Newton's identities
read
\begin{equation}\label{eq:newton-recursion}
 p_i=\sum_{j=1}^{\min(3,i)}(-1)^{j+1}s_j
 \begin{cases}
 i,&j=i,\\
 p_{i-j},&j<i.
 \end{cases}
\end{equation}
Define
\begin{equation}\label{eq:unshifted-powers}
 q_i=\sum_{j=0}^i\binom ijp_j(-h)^{i-j}.
\end{equation}
These represent $i!\operatorname{ch}_i(E)$, with $q_0=3$.
The two tangent exact sequences on $U$ give, for every $i\ge1$,
\[
 q_i+x_i=c^i+\pi^*(i!\operatorname{ch}_i(T_V)).
\]
Applying $\delta=\pi^*\pi_*$ and using~\eqref{eq:p1-powers} gives
$\delta x_i=D_i$, where
\begin{equation}\label{eq:D-polynomials}
\begin{split}
 D_i={}&(1-(-1)^i)(2h-\ell)^{i-1}\\
 &-\sum_{j=0}^{i-1}\binom ijp_j(-1)^{i-j}
      \sum_{k=0}^{i-j-1}h^k(\ell-h)^{i-j-1-k}.
\end{split}
\end{equation}
In particular $D_1=5$, consistently with $x_1=5h$, and $D_i$ evaluates
to zero for $i\ge7$. This calculation uses the polynomial product
rule for $\delta$, so remains valid even when the evaluated classes
are zero divisors.

For a monomial $m$ in the six base generators
$(h,x_2,x_3,x_4,x_5,x_6)$ define $\Delta m$ recursively by
\begin{equation}\label{eq:Delta-recursion}
 \Delta1=0,\qquad\Delta h=1,\qquad\Delta x_i=D_i\ (2\le i\le6),
 \qquad
 \Delta(gn)=(\Delta g)n+(g-c\Delta g)\Delta n.
\end{equation}
In the last formula, remove the first generator $g$ dividing $m=gn$
in the displayed order. By~\eqref{eq:p1-product}, this recursively
specified polynomial evaluates to $\delta m$. No independence of
different formal recursion orders is needed.

Finally, write $\theta_i$ for a polynomial representing
$\operatorname{td}_i(E)$. Through weight nine, put
\begin{equation}\label{eq:todd-recursion}
 \theta_0=1,\qquad
 \theta_i=\frac1i\sum_{j=1}^i\beta_jq_j\theta_{i-j},\qquad
 (\beta_1,\beta_2,\beta_4,\beta_6,\beta_8)
 =\left(\frac12,-\frac1{12},\frac1{720},-\frac1{30240},
         \frac1{1209600}\right),
\end{equation}
with the remaining $\beta_j$ zero. To obtain this formula, apply the
logarithmic derivative to the product of the Todd factors:
\[
 z\frac{d}{dz}\log\frac{z}{1-e^{-z}}
  =1-\frac{z}{e^z-1}
  =\frac z2-\frac{z^2}{12}+\frac{z^4}{720}
       -\frac{z^6}{30240}+\frac{z^8}{1209600}+O(z^{10}).
\]
Since $\operatorname{ch}(E^\vee)=\sum_{j\ge0}(-1)^jq_j/j!$, the second
polynomial in~\eqref{eq:grr-vanishings} is represented by
\begin{equation}\label{eq:todd-dual}
 \psi_i=\sum_{j=0}^i\frac{(-1)^j}{j!}q_j\theta_{i-j}.
\end{equation}

\subsection{The linear constraints and their geometric meaning}

Let $\mathcal B_k$ be the weight-$k$ monomials in
$(h,x_2,x_3,x_4,x_5,x_6)$, let $\mathcal I_k$ be those in
$(\ell,r,s,t)$, and let $\mathcal T_k$ be all weight-$k$ monomials in
$R$. A negative subscript means the empty list. Monomials are ordered
recursively: for an ordered generator list ending in a variable $g$ of
weight $a$, list $g^j$ times the monomials of weight $k-aj$ in the
preceding variables, in increasing order of $j$.

Form a rational matrix $A$ whose columns are indexed by $\mathcal T_9$.
Its rows are coefficient vectors of the following polynomials. We
group them according to their geometric source.

First, for a base monomial $m$, both $m-h\delta m$ and $\delta m$
descend to $V$ by Lemma~\ref{lem:p1-involution}, as does $\ell$.
The classes $r,s,t$ descend by Lemma~\ref{lem:shifted-descent}.
Products of these classes of
weight nine vanish because $\dim V=8$. This gives the rows
\begin{alignat}{3}
 &w(m-h\Delta m),&&\quad m\in\mathcal B_i,\quad
 w\in\mathcal I_{9-i},&&\quad0\le i\le6;\label{eq:matrix-rows-a}\\
 &w\Delta m,&&\quad m\in\mathcal B_i,\quad
 w\in\mathcal I_{10-i},&&\quad0\le i\le6.\label{eq:matrix-rows-b}
\end{alignat}
Second, a base monomial $m$ of weight greater than six vanishes,
as does $\delta m$. Also $x_i=0$ for $i\ge7$, so
$D_i=\delta x_i$ evaluates to zero. Their multiples give
\begin{alignat}{3}
 &wm,&&\quad m\in\mathcal B_i,\quad
 w\in\mathcal T_{9-i},&&\quad7\le i\le10;\label{eq:matrix-rows-c}\\
 &w\Delta m,&&\quad m\in\mathcal B_i,\quad
 w\in\mathcal T_{10-i},&&\quad7\le i\le10;\label{eq:matrix-rows-d}\\
 &wD_i,&&\quad w\in\mathcal T_{10-i},&&\quad7\le i\le10.\label{eq:matrix-rows-e}
\end{alignat}
Finally, the Riemann--Roch identities~\eqref{eq:grr-vanishings} and
the projection formula give zero integrals for
\begin{alignat}{3}
 &w\theta_i,\quad w\psi_i,&&\quad w\in\mathcal B_{9-i},&&\quad4\le i\le9.
 \label{eq:matrix-rows-f}
\end{alignat}
Zero rows may be omitted. Let
\[
 v=\left(\int_U m\right)_{m\in\mathcal T_9}.
\]
The preceding vanishings prove $Av=0$. The first five row families
come from zero cohomology classes; the final family only requires
the stated integrals to vanish.

The constraints are linear in the intersection numbers on $U$. To
recover a useful condition on $u$, we next use the product structure
forced by integration along $e$.

We use invariant monomials of weights three and four because their
pushforwards along $e$ lie in $H^0(X,\mathbb Q)$ and
$H^2(X,\mathbb Q)=\mathbb Qh_X$, respectively. Each pushforward
is determined by one scalar. These two cohomology groups
therefore supply scalar product identities.
Write the ordered lists $M=\mathcal B_6$ and
$G=(\mathcal I_3,\mathcal I_4)$ explicitly as
\begin{align*}
 M={}&(h^6,h^4x_2,h^2x_2^2,x_2^3,h^3x_3,hx_2x_3,x_3^2,
                 h^2x_4,x_2x_4,hx_5,x_6),\\
 G={}&(\ell^3,\ell^2r,\ell r^2,r^3,\ell s,rs,t,\ell^4,
       \ell^3r,\ell^2r^2,\ell r^3,r^4,\ell^2s,\ell rs,r^2s,s^2,
       \ell t,rt).
\end{align*}
Indices begin at zero, and $w_j$ denotes the weight of $G_j$.
Put $D=\int_F c_1(L)^3>0$ and define
\begin{equation}\label{eq:normalized-products}
 \xi_i=\frac{\int_X M_i}{\int_Xh_X^6},\qquad
 e_*G_j=D\eta_jh_X^{w_j-3}.
\end{equation}
The first expression uses the corresponding base classes on $X$.
The second defines a scalar because $w_j-3$ is zero or one and
$H^2(X,\mathbb Q)=\mathbb Qh_X$. We have
$\xi_0=\eta_0=1$ and $\eta_1=u$.

Use the index set
\[
 J=\{0,1,2,4,5,6,7,8,9,10,12,13,15,16,17\}
\]
for the selected invariant monomials. This subcollection suffices for
the elimination in Section~\ref{sec:arithmetic} and reduces its size;
no vanishing is asserted for the omitted monomials $r^3$, $r^4$,
and $r^2s$. Let
\[
 C_1=\{(i,(j)):0\le i\le10,\ j\in J,
                  \ w_j=3\ \text{or}\ h\mid M_i\}.
\]
Order $C_1$ by increasing $i$ and then increasing $j$.
For each such pair the monomial $M_i/h^{w_j-3}$ is defined, and the
projection formula gives
\begin{equation}\label{eq:product-factorization}
 \int_U G_j\frac{M_i}{h^{w_j-3}}
   =D\int_Xh_X^6\,\xi_i\eta_j.
\end{equation}
The common factor on the right is $\int_U h^6\ell^3>0$.
We now retain all linear relations that $Av=0$ induces on these
selected coordinates. The following construction expresses them as
a matrix $N$.
If $K$ is any matrix whose columns form a rational basis of $\ker A$,
let $P$ consist of its rows indexed, with repetitions if necessary, by
the monomials on the left of~\eqref{eq:product-factorization}.
Let $N$ be the nonzero rows in the reduced row echelon form of a matrix
whose rows form a basis of $\ker(P^{\mathsf T})$. Thus $N$ is determined
by the ordered coordinates, independently of the chosen basis $K$.
Then
\begin{equation}\label{eq:N-products}
 N(\xi_i\eta_j)_{(i,(j))\in C_1}=0.
\end{equation}
Indeed $v\in\ker A$, so its selected coordinates lie in the column
space of $P$, and division by the single positive scalar in
\eqref{eq:product-factorization} gives the displayed equation.
This completes the geometric derivation of the finite constraints.
